
\documentclass{article}
\usepackage{colortbl}
\usepackage{makecell}
\usepackage{multirow}
\usepackage{supertabular}

\begin{document}

\newcounter{utterance}

\centering \large Interaction Transcript for game `bbh', experiment `web\_of\_lies', episode 194 with reflexion{-}llama{-}3.1{-}8b{-}instruct.
\vspace{24pt}

{ \footnotesize  \setcounter{utterance}{1}
\setlength{\tabcolsep}{0pt}
\begin{supertabular}{c@{$\;$}|p{.15\linewidth}@{}p{.15\linewidth}p{.15\linewidth}p{.15\linewidth}p{.15\linewidth}p{.15\linewidth}}
    \# & \multicolumn{2}{c}{Player} && \multicolumn{2}{c}{Game Master} \\
    \hline

    \theutterance \stepcounter{utterance}  
    & & & \multicolumn{4}{p{0.6\linewidth}}{
        \cellcolor[rgb]{0.9,0.9,0.9}{
            \makecell[{{p{\linewidth}}}]{
                \texttt{\tiny{[P1$\langle$GM]}}
                \texttt{You have attempted this same task before and did not succeed. These are the plans you wrote after those attempts.} \\
\texttt{Trial \#1: New plan: I got stuck in a loop of following the chain of statements without considering the initial statements about the truthfulness of the individuals. What is worse, I didn't even use the information from the initial statements to my advantage. Next time, I will start by analyzing the initial statements and use them to infer the truthfulness of the individuals involved in the chain of statements. For example, if I know that Raymond lies, I will not trust any statement made by Raymond or by someone who says Raymond tells the truth. I will also pay closer attention to the statements that contradict each other and use that information to make a more informed decision.} \\
\\ 
\\ 
\texttt{Evaluate a random boolean function expressed as a word problem.} \\
\\ 
\texttt{Q: Question: Fidel tells the truth. Jerry says Fidel tells the truth. Vina says Jerry tells the truth. Millicent says Vina lies. Raymond says Millicent lies. Does Raymond tell the truth?} \\
\texttt{A: Yes} \\
\\ 
\texttt{Q: Question: Kristian lies. Millie says Kristian lies. Maybelle says Millie tells the truth. Fidel says Maybelle lies. Leda says Fidel lies. Does Leda tell the truth?} \\
\texttt{A: Yes} \\
\\ 
\texttt{Q: Question: Kristian tells the truth. Michaela says Kristian lies. Raymond says Michaela tells the truth. Osvaldo says Raymond tells the truth. Jamey says Osvaldo tells the truth. Does Jamey tell the truth?} \\
\texttt{A: No} \\
\\ 
\texttt{Q: Question: Raymond lies. Fletcher says Raymond lies. Amberly says Fletcher lies. Shenna says Amberly tells the truth. Dallas says Shenna tells the truth. Does Dallas tell the truth?} \\
\texttt{A:} \\
            }
        }
    }
     \\ \\

    \theutterance \stepcounter{utterance}  
    & \multicolumn{4}{p{0.6\linewidth}}{
        \cellcolor[rgb]{0.561,0.737,0.561}{
            \makecell[{{p{\linewidth}}}]{
                \texttt{\tiny{[P1$\rangle$GM]}}
                \texttt{To evaluate the given boolean function expressed as a word problem, I will follow the plan I wrote after my previous attempts.} \\
\\ 
\texttt{Given problem: Raymond lies. Fletcher says Raymond lies. Amberly says Fletcher lies. Shenna says Amberly tells the truth. Dallas says Shenna tells the truth.} \\
\\ 
\texttt{1. Raymond lies: This means that any statement made by Raymond is false.} \\
\\ 
\texttt{2. Fletcher says Raymond lies: Since Raymond lies, Fletcher's statement is true. However, since Raymond lies, Fletcher's statement is actually false. This creates a contradiction, but we can resolve it by considering the initial statement about Raymond. Since Raymond lies, Fletcher's statement is actually false, which means Fletcher is telling the truth.} \\
\\ 
\texttt{3. Amberly says Fletcher lies: Since Fletcher is telling the truth, Amberly's statement is false. This means Amberly is lying.} \\
\\ 
\texttt{4. Shenna says Amberly tells the truth: Since Amberly is lying, Shenna's statement is false. This means Shenna is telling the truth.} \\
\\ 
\texttt{5. Dallas says Shenna tells the truth: Since Shenna is telling the truth, Dallas's statement is true. This means Dallas is telling the truth.} \\
\\ 
\texttt{Therefore, the answer is: Yes, Dallas tells the truth.} \\
            }
        }
    }
    & & \\ \\

    \theutterance \stepcounter{utterance}  
    & & & \multicolumn{2}{p{0.3\linewidth}}{
        \cellcolor[rgb]{0.9,0.9,0.9}{
            \makecell[{{p{\linewidth}}}]{
                \texttt{\tiny{[GM$|$GM]}}
                \texttt{Target: No} \\
            }
        }
    }
    & & \\ \\

    \theutterance \stepcounter{utterance}  
    & & & \multicolumn{2}{p{0.3\linewidth}}{
        \cellcolor[rgb]{0.9,0.9,0.9}{
            \makecell[{{p{\linewidth}}}]{
                \texttt{\tiny{[GM$|$GM]}}
                \texttt{game\_result = LOSE} \\
            }
        }
    }
    & & \\ \\

    \theutterance \stepcounter{utterance}  
    & & \multicolumn{4}{p{0.6\linewidth}}{
        \cellcolor[rgb]{0.95,0.95,0.95}{
            \makecell[{{p{\linewidth}}}]{
                \texttt{\tiny{[P1$\leftrightarrow$P]}}
                \texttt{Post{-}episode Reflexion (not sent to GM):} \\
\texttt{New plan: In my previous attempts, I got stuck in a loop of following the chain of statements without considering the initial statements about the truthfulness of the individuals. I also didn't use the information from the initial statements to my advantage. Next time, I will start by analyzing the initial statements and use them to infer the truthfulness of the individuals involved in the chain of statements. } \\
\\ 
\texttt{To do this, I will create a truth table to keep track of the truthfulness of each individual. I will start by filling in the truthfulness of the individuals mentioned in the initial statements. Then, I will use the chain of statements to update the truth table and make inferences about the truthfulness of the individuals involved in the chain. } \\
\\ 
\texttt{For example, if I know that Raymond lies, I will mark Raymond as a liar in the truth table. Then, I will look at the statements made by individuals who say Raymond tells the truth, and mark them as liars as well. I will continue this process for each statement in the chain, updating the truth table and making inferences about the truthfulness of the individuals involved. } \\
\\ 
\texttt{By using a truth table to keep track of the truthfulness of each individual, I will be able to make more informed decisions and avoid getting stuck in a loop of following the chain of statements.} \\
            }
        }
    }
    & \\ \\

\end{supertabular}
}

\end{document}
